\documentclass[12pt]{article}
\usepackage{amsmath, amssymb, amsthm}
\usepackage{geometry}
\geometry{margin=1in}

\newtheorem{exercise}{Exercise}
\newtheorem{solution}{Solution}

\title{Random Walks --- Week 1 Piazza Exercises}
\author{CUNY, Spring 2026}
\date{}

\begin{document}
\maketitle

\section*{Week 1: Foundations of Simple Random Walk}

Recall the finite-time setup: $\Omega_N = \{-1,+1\}^N$ with the uniform probability measure
\[
  P_N(A) = \frac{|A|}{2^N}, \qquad A \subseteq \Omega_N.
\]
Define $X_k(\omega) = \omega_k$ for $\omega = (\omega_1, \ldots, \omega_N) \in \Omega_N$ and $k = 1, \ldots, N$.

\begin{exercise}
Show that $X_1, X_2, \ldots, X_N$ are independent and identically distributed with
\[
  P_N(X_i = +1) = P_N(X_i = -1) = \tfrac{1}{2}.
\]
\end{exercise}

\begin{proof}
\textbf{Identical distribution.} Fix any $k \in \{1,\ldots,N\}$. The set
\[
  \{X_k = +1\} = \{\omega \in \Omega_N : \omega_k = +1\}
\]
has exactly $2^{N-1}$ elements (we freely choose each of the remaining $N-1$ coordinates). Therefore
\[
  P_N(X_k = +1) = \frac{2^{N-1}}{2^N} = \frac{1}{2},
\]
and by an identical argument $P_N(X_k = -1) = \tfrac{1}{2}$. Since this holds for every $k$, the $X_k$ are identically distributed.

\textbf{Independence.} Let $k_1 < k_2 < \cdots < k_n$ be distinct indices in $\{1,\ldots,N\}$, and let $x_1, \ldots, x_n \in \{-1,+1\}$. The event
\[
  \{X_{k_1} = x_1\} \cap \cdots \cap \{X_{k_n} = x_n\}
\]
fixes $n$ specific coordinates of $\omega$ and leaves the remaining $N-n$ coordinates free, so it contains exactly $2^{N-n}$ elements. Hence
\[
  P_N\!\bigl(X_{k_1}=x_1,\,\ldots,\,X_{k_n}=x_n\bigr)
  = \frac{2^{N-n}}{2^N} = \frac{1}{2^n}
  = \prod_{i=1}^n P_N(X_{k_i} = x_i).
\]
Since this factorization holds for all choices of indices and values, $X_1, \ldots, X_N$ are mutually independent.
\end{proof}

\end{document}
